\section{La cérémonie du jonc}\label{la-ceremonie-du-jonc}

\stanceinfo{\textbf{CESS de la FCÉG 2023} [2023-03-02]}{2024-01-02}

\officialstance{La Fédération canadienne étudiante de génie estime que la cérémonie du jonc est une cérémonie de transition nécessaire pour les étudiants en génie, mais que, dans son état actuel, elle est obsolète et ne représente plus pleinement l'ingénierie de l'ère moderne. Elle limite plutôt la description de la profession d'ingénieur et possède un ton d'exclusion, et elle devrait être modifiée pour reconnaître l'obligation de tous les ingénieurs.}

\stanceheading{La position des étudiants}
\begin{itemize}
 \item La cérémonie du jonc est un rite de passage essentiel pour les étudiants en ingénierie. Elle est très appréciée des diplômés après l'obtention d'un diplôme long et difficile.
 \item Dans le meilleur des cas, la cérémonie du jonc peut contribuer à un sentiment d'appartenance et de communauté dans le domaine de l'ingénierie.
 \item Le rituel original ne s'adresse pas à tous les étudiants en ingénierie, ce qui les déçoit. Le rituel restreint la pratique de l'ingénierie en dehors des formes de travail originales lorsque le rituel a été promulgué pour la première fois et présente une vision du monde dépassée en ce qui concerne le colonialisme, le racisme et le sexisme.
 \item Chaque camp organise sa propre cérémonie avec peu de surveillance de la part de la corporation des sept gardiens.
 \item Tous les étudiants en ingénierie devraient se sentir inclus, à l'aise et respectés pendant le Rituel de l'appel de l'ingénieur.
\end{itemize}

\stanceheading{L'enjeu}

La Cérémonie du jonc est l'aboutissement d'années de dur travail pour les étudiants en ingénierie au Canada, un rituel au cours duquel l'étudiant est appelé à prendre l'engagement de pratiquer l'ingénierie selon des normes éthiques et professionnelles [7]. C'est pourquoi la cérémonie revêt une grande importance pour les étudiants de premier cycle. La fonction principale de la cérémonie est de faire en sorte que les étudiants en ingénierie se sentent positifs et affirmés en tant qu'ingénieurs.

Présentement, certains étudiants de premier cycle quittent la cérémonie avec un sentiment de manque d'accomplissement et d'appartenance à la profession d'ingénieur, s'éloignant ainsi de l'objectif initial. Outre les voix d'étudiants exprimant leur malaise et leur mécontentement à l'égard de l'actuel Rituel d'appel de l'ingénieur, des ingénieurs et des formateurs en ingénierie se sont unis pour rédiger une déclaration soumise au Conseil des sept gardiens en 2022, avec un mouvement intitulé « Retrouver l'anneau », soulignant principalement qu'il n'incarne pas une compréhension globale de l'éthique de l'ingénierie, qu'il manque de clarté et de transparence dans le texte et qu'il véhicule des visions du monde dépassées et nuisibles, enracinées dans le colonialisme, le racisme et le sexisme [8]. Ils ont demandé au Conseil de reconnaître et d'assumer ses responsabilités afin de créer un nouvel appel de l'ingénieur qui englobe la pratique moderne de l'ingénierie.

Le plus grand défi de la Cérémonie du jonc réside dans le fait que ces questions ne se posent pas uniformément, car il existe actuellement plusieurs formes du poème « Le rite de l'appel de l'ingénieur ». Chaque camp organise la Cérémonie du jonc de manière indépendante pour une certaine zone géographique du pays. Chaque camp est autorisé à créer sa propre version de la cérémonie, tandis que d'autres ont résolu les problèmes susmentionnés. Certains camps ont ouvert la cérémonie aux invités, d'autres ont commencé à commander de nouveaux écrits à utiliser lors de la cérémonie.

La FCEG est convaincue que tous les étudiants en ingénierie devraient être encouragés à devenir ingénieurs lors du rite de ``l'appel de l'ingénieur'' et qu'aucun étudiant ne devrait se sentir irrespectueux, mal à l'aise ou diminué en raison de son sexe, de son âge, de sa religion, de sa race ou de tout autre attribut extérieur qui n'affecte pas la pratique de l'ingénierie. En tant que moment clé, ce rite doit éclairer et encourager les ingénieurs de demain à pratiquer et à se développer dans leur profession d'ingénieur.

\stanceheading{Les recommandations aux partenaires, aux parties intéressées et aux autres entités}
\begin{itemize}
 \item Les partenaires, les intervenants et les autres entités sont encouragés à collaborer avec la FCEG pour veiller à ce que les recommandations des membres soient prises en compte lors de l'exécution de la Cérémonie du jonc.
 \item Intégrer les voix des étudiants de premier et de deuxième cycle dans le comité actuel intitulé « Comité de révision du rituel » [9] pour la révision de la Cérémonie du jonc, sur la base du Conseil des sept gardiens.
 \item Que les membres recueillent directement auprès des diplômés des informations sur les changements qu'ils souhaiteraient voir se produire, soit par l'intermédiaire des organisations régionales, soit par celui de leur propre société d'ingénieurs.
 \item La FCEG demande à Ingénieurs Canada de faire pression pour que des informations soient données aux étudiants de premier cycle avant qu'ils n'assistent à la cérémonie
\end{itemize}

\stanceheading{Ce que fait la FCEG}
\begin{itemize}
 \item Collaborer avec les organisations membres pour partager les connaissances et la compréhension de la Cérémonie du jonc
 \item Recueillir les réactions des diplômés qui ont participé à la Cérémonie du jonc.
\end{itemize}

\stanceheading{Ce que la FCEG envisage de faire}
\begin{itemize}
 \item Veiller à ce que la Cérémonie du jonc soit correctement supervisée par le Conseil des sept gardiens, afin que tous les élèves-ingénieurs puissent vivre une expérience positive lors de la Cérémonie du jonc.
 \item Réaliser une enquête auprès des élèves ingénieurs de tous les camps sur leur expérience de la Cérémonie du jonc.
 \item Ateliers/présentations visant à diffuser des informations essentielles
 \item Enquête sur la Cérémonie du jonc auprès des étudiants actuels et des jeunes diplômés.
 \item Compilation d'enquêtes préexistantes
 \item La FCEG s'engage à organiser au moins cinq présentations ou ateliers éducatifs par an, visant à diffuser des informations essentielles sur la Cérémonie du jonc. Ces sessions s'adresseront à nos membres, afin de garantir une large compréhension de la cérémonie.
 \item La FCEG réalisera une enquête détaillée (date à déterminer) dans le but de recueillir des informations auprès d'un nombre minimum de participants (à déterminer), comprenant des étudiants actuels et des diplômés récents. Cette enquête portera spécifiquement sur les thèmes de l'inclusion, de la représentation, de l'impact et de la logistique, en veillant à ce que le retour d'information soit pertinent, généralisé et exploitable. Les résultats de cette enquête seront examinés chaque année afin de suivre les progrès accomplis et l'évolution de la perception des étudiants.
\end{itemize}

\newpage
